\documentclass[11pt,a4paper]{article}
\usepackage[utf8]{inputenc}
\usepackage[vietnamese]{babel}
\usepackage{amsmath,amssymb}
\usepackage{geometry}
\geometry{a4paper, total={180mm,260mm}, left=15mm, right=15mm, top=18mm, bottom=18mm}
\usepackage{tikz}
\usetikzlibrary{shapes.geometric, arrows.meta, positioning, calc}
\usepackage{tcolorbox}
\tcbuselibrary{skins,breakable}
\usepackage{xcolor}
\usepackage{multicol}
\usepackage{array}
\usepackage{fontawesome}
\usepackage{fancyhdr}

\pagestyle{fancy}
\fancyhf{}
\renewcommand{\headrulewidth}{0.5pt}
\renewcommand{\footrulewidth}{0.5pt}
\fancyhead[L]{\small\bfseries\color{teal}\faBook\ SỔ TAY KIẾN THỨC TOÁN 3}
\fancyhead[R]{\small\bfseries\color{teal}Lớp Toán Thầy Lĩnh -- loptoanthaylinh.vn}
\fancyfoot[C]{\small\color{gray}-- Sổ tay kiến thức: Bài 29. Luyện tập chung: Phép nhân và phép chia phạm vi 100 --}

\definecolor{primary}{HTML}{005F73}
\definecolor{secondary}{HTML}{0A9396}
\definecolor{accent}{HTML}{EE9B00}
\definecolor{danger}{HTML}{AE2012}
\definecolor{boxbg}{HTML}{F4F6F8}

\begin{document}

% --- BANNER TIÊU ĐỀ ---
\begin{tcolorbox}[
  enhanced,
  colback=white,
  colframe=primary,
  arc=3mm,
  boxrule=1.2pt,
  left=4mm, right=4mm, top=2.5mm, bottom=2.5mm,
  drop shadow=black!10!white
]
  \centering
  {\Large\bfseries\color{primary}\faBookmark\ BÀI 29: LUYỆN TẬP CHUNG: NHÂN CHIA PHẠM VI 100}\\[1.5mm]
  {\small\itshape\color{black!75}Đại tổng kết đặt tính, tính giá trị biểu thức, gấp giảm số lần và bài toán hai bước trên 1 trang A4}
\end{tcolorbox}

\vspace{-2mm}


\noindent
\begin{minipage}[t]{0.49\textwidth}\vspace{0pt}
\begin{tcolorbox}[
  enhanced,
  colback=blue!4!white,
  colframe=blue!75!black,
  arc=2mm, boxrule=0.8pt,
  left=2mm, right=2mm, top=1.8mm, bottom=1.8mm,
  title={\footnotesize\bfseries\color{white}\faCalculator\ 1. Kĩ thuật tính toán 4 quy tắc}
]
{\footnotesize
\begin{itemize}\setlength{\itemsep}{1.5pt}\setlength{\leftskip}{-2.5mm}
  \item \textbf{Phép nhân $2$ chữ số với $1$ chữ số:}
  Nhân lần lượt từ phải sang trái (nhớ sang hàng chục nếu có).
  \item \textbf{Phép chia $2$ chữ số cho $1$ chữ số:}
  Chia lần lượt từ trái sang phải qua 3 nhịp: Chia $\to$ Nhân $\to$ Trừ.
  \item \textbf{Quy tắc số dư:} Số dư luôn \textbf{nhỏ hơn} số chia:
  $$r < d \implies \text{Số bị chia} = \text{Thương} \times \text{Số chia} + r$$
\end{itemize}
}
\end{tcolorbox}
\end{minipage}
\hfill
\begin{minipage}[t]{0.49\textwidth}\vspace{0pt}
\begin{tcolorbox}[
  enhanced,
  colback=teal!4!white,
  colframe=teal!75!black,
  arc=2mm, boxrule=0.8pt,
  left=2mm, right=2mm, top=1.8mm, bottom=1.8mm,
  title={\footnotesize\bfseries\color{white}\faCheckSquare\ 2. Kĩ năng giải toán tổng hợp}
]
{\footnotesize
\begin{itemize}\setlength{\itemsep}{1.5pt}\setlength{\leftskip}{-2.5mm}
  \item \textbf{Gấp $n$ lần:} Làm phép nhân $(\times n)$.
  \item \textbf{Giảm $n$ lần:} Làm phép chia $(: n)$.
  \item \textbf{Tìm một phần mấy của một số:}
  Muốn tìm $\dfrac{1}{n}$ của số $A$, ta lấy $A : n$.
  \item \textbf{Thứ tự thực hiện biểu thức:}
  $$\text{Trong ngoặc trước} \to \text{Nhân/Chia} \to \text{Cộng/Trừ}$$
\end{itemize}
}
\end{tcolorbox}
\end{minipage}

\vspace{-2.5mm}
% --- KHỐI 3: SƠ ĐỒ TƯ DUY TIKZ ---
\noindent
\begin{tcolorbox}[
  enhanced,
  colback=orange!4!white,
  colframe=orange!80!black,
  arc=2mm, boxrule=0.8pt,
  left=2mm, right=2mm, top=1.5mm, bottom=1.5mm,
  title={\footnotesize\bfseries\color{white}\faSitemap\ SƠ ĐỒ TƯ DUY: HỆ THỐNG 4 NỀN TẢNG NHÂN CHIA LỚP 3}
]
\centering
\begin{tikzpicture}[
  box/.style={draw, rounded corners=2pt, align=center, font=\scriptsize\bfseries, inner sep=2.5pt, minimum height=0.75cm},
  arr/.style={-{Stealth[length=2mm]}, thick}
]
  \node[box, fill=purple!15, draw=purple!80!black, text width=2.6cm] (d1) at (0, 0) {BẢNG CỬU CHƯƠNG\\Nhân chia 2 - 9\\Mối quan hệ thuận nghịch};
  \node[box, fill=teal!15, draw=teal!80!black, text width=2.6cm] (d2) at (3.5, 0) {ĐẶT TÍNH\\Nhân phải qua trái\\Chia trái qua phải};
  \node[box, fill=blue!12, draw=blue!80!black, text width=2.6cm] (d3) at (7.0, 0) {GẤP / GIẢM\\Gấp $\to$ Nhân $\times$\\Giảm $\to$ Chia $:$};
  \node[box, fill=orange!15, draw=orange!80!black, text width=2.6cm] (d4) at (10.5, 0) {TOÁN 2 BƯỚC\\Tìm ẩn trung gian\\Tính tổng hoặc còn lại};

  \draw[arr, purple!80!black] (d1) -- (d2);
  \draw[arr, teal!80!black] (d2) -- (d3);
  \draw[arr, blue!80!black] (d3) -- (d4);
\end{tikzpicture}
\end{tcolorbox}

\vspace{-2.5mm}
% --- KHỐI 4: VÍ DỤ & MẸO TRÁNH BẪY ---
\noindent
\begin{minipage}[t]{0.49\textwidth}\vspace{0pt}
\begin{tcolorbox}[
  enhanced,
  colback=green!4!white,
  colframe=green!60!black,
  arc=2mm, boxrule=0.8pt,
  left=2mm, right=2mm, top=1.8mm, bottom=1.8mm,
  title={\footnotesize\bfseries\color{white}\faCheckCircle\ Ví dụ thực hành chuẩn mực}
]
{\footnotesize
\textbf{1. Đặt tính chia có dư:}\\
$75 : 6 = 12$ (dư $3$).\\
Thử lại: $12 \times 6 + 3 = 72 + 3 = 75$ (đúng).\\
\textbf{2. Tính nhanh giá trị biểu thức:}\\
$A = 24 \times 3 - 35 : 5 = 72 - 7 = \mathbf{65}$.\\
$B = (15 + 25) : 4 = 40 : 4 = \mathbf{10}$.
}
\end{tcolorbox}
\end{minipage}
\hfill
\begin{minipage}[t]{0.49\textwidth}\vspace{0pt}
\begin{tcolorbox}[
  enhanced,
  colback=red!4!white,
  colframe=red!70!black,
  arc=2mm, boxrule=0.8pt,
  left=2mm, right=2mm, top=1.8mm, bottom=1.8mm,
  title={\footnotesize\bfseries\color{white}\faExclamationTriangle\ Mẹo tránh bẫy trắc nghiệm}
]
{\footnotesize
\begin{itemize}\setlength{\itemsep}{1.5pt}\setlength{\leftskip}{-2.5mm}
  \item \textbf{Bẫy tính từ trái sang phải mà quên thứ tự:}\\
  Ví dụ biểu thức $20 + 30 : 5$, nhiều bạn tính $20 + 30 = 50$ rồi chia cho $5$ bằng $10$ là \textbf{SAI}! Đúng phải là $30 : 5 = 6 \implies 20 + 6 = \mathbf{26}$.
  \item \textbf{Bẫy quên số dư lớn hơn số chia:}\\
  Nếu số dư $\ge$ số chia thì thương tìm được chưa tối ưu, phải tăng thương lên!
\end{itemize}
}
\end{tcolorbox}
\end{minipage}


\end{document}
